\documentclass[12pt]{book}
\usepackage{inputenc}
\usepackage[spanish]{babel}
\decimalpoint % Usar puntos (.) en lugar de comas (,) como separadores decimales
\usepackage{geometry}
\usepackage{fancyhdr}
\usepackage{graphicx}
\usepackage{amsmath} % Para entornos matemáticos avanzados
\usepackage{amssymb} % Para el uso de \mathbb
\usepackage{amsthm} % Para entornos como definition, proposition, example, remark
\usepackage{hyperref} % Para hipervínculos interactivos en el PDF

\geometry{margin=1in}
\pagestyle{fancy}

% Configuración de hyperref para mejor apariencia
\hypersetup{
    colorlinks=true,
    linkcolor=blue,
    citecolor=blue,
    urlcolor=blue,
    bookmarksnumbered=true,
    pdfstartview=FitH
}

% Ocultar el texto de las secciones en los encabezados
\renewcommand{\sectionmark}[1]{}

% Definición de entornos personalizados
\newtheorem{theorem}{Teorema}[chapter]
\newtheorem{definition}{Definición}[chapter]
\newtheorem{proposition}{Proposición}[chapter]
\newtheorem{example}{Ejemplo}[chapter]
\newtheorem{remark}{Observación}[chapter]

% Ajuste de \headheight para evitar advertencias de fancyhdr
\setlength{\headheight}{15.71667pt}

\title{Procesos Estocásticos I}
\author{Sonny Medina}
\date{Semestre 2026-2}

\begin{document}

% Portada
\begin{titlepage}
    \centering
    \vspace*{1cm}
    
    {\Large \textbf{UNIVERSIDAD NACIONAL AUTÓNOMA DE MÉXICO}}\\
    \vspace{0.3cm}
    {\large Facultad de Estudios Superiores Acatlán}\\
    
    \vspace{3cm}
    
    {\Huge \textbf{Procesos Estocásticos I}}\\
    
    \vspace{2cm}
    
    {\Large Semestre 2026-2}\\
    
    \vspace{3cm}
    
%    {\large \textbf{Profesor:}}\\
    {\large Sonny Medina (Draft)}\\
    
    \vfill
    
    {\large \today}
    
\end{titlepage}

% Tabla de contenidos
\tableofcontents
\newpage

% Estructura de capítulos
\chapter{Modelos Markovianos en Tiempo Discreto}

%\section*{Introducción Histórica}
La teoría de procesos estocásticos tuvo sus orígenes en los trabajos de Andréi Márkov a principios del siglo XX. Márkov, discípulo de Pafnuty Chebyshev al igual que Andréi Kolmogorov, estudió sucesiones de variables aleatorias que satisfacían una propiedad especial: la dependencia entre observaciones sucesivas. Aunque las contribuciones de Márkov no fueron tan formales como las posteriores de Kolmogorov, quien desarrolló la teoría axiomática de la probabilidad, Márkov es considerado el precursor de la teoría moderna de procesos estocásticos. En sus trabajos iniciales, Márkov se enfocó principalmente en procesos a tiempo discreto con un enfoque probabilístico. Con el desarrollo de la teoría, especialmente tras las contribuciones de Kolmogorov, la disciplina se formalizó considerablemente e incorporó tanto procesos a tiempo discreto como continuo, ampliando su alcance hacia aplicaciones contemporáneas en física estadística, finanzas matemáticas, biología, ingeniería y ciencias de la computación. En la actualidad, los procesos estocásticos constituyen una herramienta fundamental para modelar sistemas dinámicos bajo incertidumbre, con aplicaciones que van desde la modelación de mercados financieros y sistemas de colas, hasta algoritmos de aprendizaje automático, métodos de Monte Carlo vía cadenas de Markov y el análisis de datos temporales en contextos científicos e industriales.

\section{Definición de Proceso Estocástico}

\begin{definition} Un \textbf{proceso estocástico} es una colección de variables aleatorias $\{X_t : t \in T\}$ definidas sobre un espacio de probabilidad $(\Omega, \mathcal{F}, \mathbb{P})$, donde:
\begin{itemize}
    \item $T$ es el \textbf{índice temporal}, que puede ser discreto ($T = \mathbb{N}$) o continuo ($T = \mathbb{R}^+$). Asumiremos que ambos conjuntos contienen al cero.
    \item $X_t$ representa el estado del sistema en el tiempo $t$.
\end{itemize}


Un proceso estocástico es \textbf{discreto} si $T$ es un conjunto discreto, y \textbf{continuo} si $T$ es un conjunto continuo.
\end{definition}

\subsection*{Interpretación como función}

Un proceso estocástico puede interpretarse como una función
\[
X : T \times \Omega \to S
\]
tal que a cada pareja $(t, \omega) \in T \times \Omega$ se le asocia el estado $X(t, \omega) \in S$. Para cada tiempo $t \in T$ fijo, la función $\omega \mapsto X(t, \omega)$ es una variable aleatoria. Por el contrario, para cada resultado $\omega \in \Omega$ fijo, la función $t \mapsto X(t, \omega)$ se llama una \textbf{trayectoria} o \textbf{realización} del proceso, y representa la evolución temporal del sistema bajo el escenario particular $\omega$.

\vspace{10pt}
\hrule 
{\center \textit{EJEMPLOS} }

\section{Espacio de Estados}

El \textbf{espacio de estados} de un proceso estocástico es el conjunto $S$ de todos los valores posibles que las variables aleatorias $X_t$ pueden tomar. Se requiere que $S$ sea equipado con una $\sigma$-álgebra apropiada para poder hablar de eventos relacionados con los valores del proceso. Consideraremos principalmente espacios de estados discretos (subconjuntos de $\mathbb{Z}$) o continuos (subconjuntos de $\mathbb{R}$, $\mathbb{R}^d$, etc.). Por ejemplo:
\begin{itemize}
    \item Si $S = \{0, 1\}$, el proceso es binario.
    \item Si $S = \mathbb{R}$, el proceso tiene un espacio de estados continuo.
    \item Si $S = \mathbb{Z}$, el proceso tiene un espacio de estados discreto infinito.
    \item Si $S = \mathbb{R}^d$ para $d \geq 2$, el espacio de estados es multidimensional.
\end{itemize}

\begin{example}[1]
\textbf{Sucesión de resultados de un dado.} Considere un dado justo. El espacio parametral es $T = \mathbb{N}$, ya que observamos los resultados en pasos discretos (cada lanzamiento). El espacio de estados es $S = \{1, 2, 3, 4, 5, 6\}$ y supondremos que s.p.g. $P(X_0=6)=1$. Ejemplos de probabilidades de interés incluyen:
\begin{itemize}
    \item $\mathbb{P}(X_3 = 6)$: Probabilidad de obtener un 6 en el tercer lanzamiento.
    \item $\mathbb{P}( \text{al menos un 6 en 5 lanzamientos} )$.
    \item $\mathbb{P}( \text{El primer 6 sale en el lanzamiento } k )$.
    \item Probabilidades condicionales como $\mathbb{P}(X_4 = 5 \mid X_3 = 2)$.
\end{itemize}
\end{example}

Un proceso estocástico se llama cadena cuando es un proceso indexado en tiempo discreto, con espacio de estados numerable. Si dibujamos los posibles estados como nodos en un grafo, y las transiciones entre estados como aristas dirigidas obtenemos una representación visual útil para analizar el comportamiento del proceso.

\section{Caminata Aleatoria}

Un ejemplo fundamental en la teoría de procesos estocásticos es la \textbf{caminata aleatoria}. Sea $\{Y_i : i \in \mathbb{N}\}$ una sucesión de variables aleatorias independientes e idénticamente distribuidas (i.i.d.), todas definidas sobre el mismo espacio de probabilidad $(\Omega, \mathcal{F}, \mathbb{P})$, que toman valores en $S_Y \subseteq \mathbb{R}$. Denotamos por $\mu_Y$ la distribución común de las variables $Y_i$. Consideremos la sucesión de sumas parciales:
\[
X_n = \sum_{i=1}^n Y_i, \quad n \in \mathbb{N}, \quad \text{con } X_0 = 0.
\]

Este proceso se llama \textbf{caminata aleatoria} o \textbf{proceso de suma acumulada}, y es un proceso a tiempo discreto con espacio de estados $S \subseteq \mathbb{R}$ (típicamente $S = \mathbb{Z}$ o $S = \mathbb{R}$ dependiendo de la distribución de $Y_i$).

\begin{proposition}
Para una caminata aleatoria $X_n = \sum_{i=1}^n Y_i$ donde los $Y_i$ son i.i.d.:
\begin{itemize}
%    \item La distribución inicial es $\mathbb{P}(X_0 = 0) = 1$.
    \item La distribución de $X_n$ es la convolución de $n$ copias de $\mu_Y$.
    \item Para $n \geq 1$, se tiene $X_n = X_{n-1} + Y_n$, por lo que la probabilidad de transición depende únicamente de la distribución de $Y_i$.
\end{itemize}
\end{proposition}

\begin{example}[2]
\textbf{Caminata aleatoria simétrica en $\mathbb{Z}$.} Sea $Y_i$ con distribución $\mathbb{P}(Y_i = 1) = \mathbb{P}(Y_i = -1) = \frac{1}{2}$. El proceso $X_n = \sum_{i=1}^n Y_i$ es una caminata aleatoria simétrica donde en cada paso el proceso se mueve una unidad hacia la derecha o hacia la izquierda con igual probabilidad. Algunas probabilidades de interés:
\begin{itemize}
    \item $\mathbb{P}(X_1 = 1) = \frac{1}{2}$ y $\mathbb{P}(X_1 = -1) = \frac{1}{2}$.
    \item $\mathbb{P}(X_2 = 0) = \mathbb{P}(Y_1 = 1, Y_2 = -1) + \mathbb{P}(Y_1 = -1, Y_2 = 1) = \frac{1}{4} + \frac{1}{4} = \frac{1}{2}$.
    \item $\mathbb{P}(X_2 = 2) = \mathbb{P}(Y_1 = 1, Y_2 = 1) = \frac{1}{4}$.
    \item $\mathbb{P}(X_3 = 0)$: Es necesario un número par de pasos $+1$ y un número par de pasos $-1$. Esto es imposible en 3 pasos, por lo que esta probabilidad es 0.
\end{itemize}
\end{example}

\begin{example}[3]
\textbf{Caminata aleatoria de Pólya.} Consideremos $Y_i$ con distribución $\mathbb{P}(Y_i = k) = p^k(1-p)^{1-k}$ para $k \in \{0, 1\}$, donde $0 < p < 1$. El proceso $X_n = \sum_{i=1}^n Y_i$ cuenta el número de ``éxitos'' hasta el tiempo $n$, y representa el número de caras en $n$ lanzamientos de una moneda sesgada. Este tipo de caminata aleatoria tiene aplicaciones en economía, biología y otras áreas.
\end{example}

\section{Distribución Inicial y Matriz de Transición}

\textbf{Distribución inicial.} La \textbf{distribución inicial} (o \textbf{distribución de probabilidad inicial}) de un proceso estocástico describe la ley de probabilidad del estado del sistema en el tiempo inicial. Para un proceso con espacio de estados $S$, se define como:
\[
\pi = (\pi_i)_{i \in S}, \quad \text{donde} \quad \pi_i = \mathbb{P}(X_0 = i) \geq 0, \quad \sum_{i \in S} \pi_i = 1.
\]

En el caso de espacios de estados infinitos, la distribución inicial es una medida de probabilidad sobre $S$.

\textbf{Probabilidades de transición.} Para procesos a tiempo discreto, la \textbf{probabilidad de transición} de un paso describe la probabilidad de pasar de un estado $i$ al estado $j$ entre tiempo $n$ y $n+1$:
\[
p_{ij} = \mathbb{P}(X_{n+1} = j \mid X_n = i), \quad \forall i, j \in S.
\]

Cuando el espacio de estados es finito, estas probabilidades pueden organizarse en una \textbf{matriz de transición}:
\[
P = (p_{ij})_{i,j \in S},
\]
que es una matriz cuadrada cuyas entradas satisfacen propiedades específicas.

\begin{proposition}
Sea $P = (p_{ij})$ la matriz de probabilidades de transición de un paso para un proceso estocástico con espacio de estados finito $S = \{1, 2, \ldots, m\}$. Entonces:
\begin{enumerate}
    \item $p_{ij} \geq 0$ para todos $i, j \in S$.
    \item $\sum_{j=1}^m p_{ij} = 1$ para todo $i \in S$ (cada fila suma 1).
    \item $P$ es una \textbf{matriz estocástica}.
\end{enumerate}
\end{proposition}

\begin{example}[4]
\textbf{Movilidad social y clases socioeconómicas.} Considere una población clasificada en tres clases: baja (B), media (M) y alta (A). Sea $X_n$ la clase socioeconómica del primogénito de una familia en la generación $n$. El espacio de estados es $S = \{B, M, A\}$ y el espacio parametral es $T = \mathbb{N}$. 

Suponga que en el momento actual ($n = 0$), la distribución inicial es:
\[
\pi_B = 0.5, \quad \pi_M = 0.3, \quad \pi_A = 0.2.
\]

Esto significa que el 50\% de la población pertenece a la clase baja, el 30\% a la clase media y el 20\% a la clase alta. Una pregunta natural es: \textit{¿Cuál es la probabilidad de que el nieto de una persona de clase baja pertenezca a la clase alta?} Para responder esto necesitamos conocer las probabilidades de transición $p_{ij} = \mathbb{P}(X_{n+1} = j \mid X_n = i)$ para cada par de clases $(i,j)$.
\end{example}

\begin{example}[5]
\textbf{Dinero en la cartera.} Sea $X_n$ la cantidad de dinero (en pesos) en nuestra cartera al final del día $n$. El espacio de estados es $S = \mathbb{R}^+_0$ (números reales no negativos). El espacio parametral es $T = \mathbb{N}$ (días). La distribución inicial $X_0$ depende de cuánto dinero teníamos al inicio, y las probabilidades de transición modelan cómo cambia nuestro saldo diariamente en función de nuestros gastos e ingresos, que pueden ser aleatorios.
\end{example}

\begin{example}[6]
\textbf{Precio de una acción en la bolsa.} Sea $X_n$ el precio de una acción particular al cierre del día $n$. El espacio parametral es $T = \mathbb{N}$ (días de operación bursátil) y el espacio de estados es $S = \mathbb{R}^+$ (precios positivos). La distribución inicial depende del precio de cierre inicial de la acción. Las probabilidades de transición $\mathbb{P}(X_{n+1} \in [a, b] \mid X_n = x)$ modelan cómo evoluciona el precio día a día en respuesta a condiciones de mercado, y son un objeto central de estudio en matemática financiera.
\end{example}

\begin{example}[7]
\textbf{Temperatura ambiente cada hora.} Sea $X_t$ la temperatura ambiente (en grados Celsius) registrada cada hora $t$ durante un periodo de tiempo. Aquí, $T = \mathbb{N}$ (horas) y $S = \mathbb{R}$ (cualquier temperatura es posible en principio). La distribución inicial $X_0$ depende de la temperatura al momento inicial, y las probabilidades de transición reflejan patrones de cambio de temperatura entre horas consecutivas, que dependen de factores como la hora del día, la estación del año, y otras condiciones meteorológicas.
\end{example}

\begin{example}[8]
\textbf{Deciles de ingresos.} Sea $X_n$ el decil de ingresos al que pertenece el primogénito de una familia en la generación $n$, donde los deciles son $S = \{D_1, D_2, \ldots, D_{10}\}$ con $D_1$ siendo el 10\% más pobre y $D_{10}$ el 10\% más rico. El espacio parametral es $T = \mathbb{N}$ (generaciones). La distribución inicial $(\pi_{D_1}, \ldots, \pi_{D_{10}})$ representa la proporción de la población en cada decil. Las probabilidades de transición $p_{ij} = \mathbb{P}(X_{n+1} = D_j \mid X_n = D_i)$ describen la movilidad socioeconómica entre generaciones.
\end{example}

\begin{example}[9]
\textbf{Estatura de una persona en el tiempo.} Sea $X_n$ la estatura de una persona (en centímetros) medida al final del año $n$. Aquí, $T = \mathbb{N}$ (años de vida) y $S = \mathbb{R}^+$ (estaturas positivas). La distribución inicial depende de la estatura al nacimiento. Las probabilidades de transición son determinísticas en el caso del crecimiento infantil (la estatura tiende a aumentar de manera predecible años a años durante la infancia), pero pueden incluir variabilidad en períodos posteriores.
\end{example}

\begin{example}[10]
\textbf{Población de una especie.} Sea $X_n$ la cantidad total de individuos de una especie particular en una región al tiempo $n$. El espacio de estados es $S = \mathbb{Z}_{\geq 0}$ (números enteros no negativos), el espacio parametral es $T = \mathbb{N}$ (periodos de tiempo, p. ej., años), y la distribución inicial es el tamaño de la población inicial. Las probabilidades de transición dependen de tasas de natalidad, mortalidad, inmigración y emigración, que pueden ser determinísticas o estocásticas.
\end{example}

\begin{example}[11]
\textbf{Movimiento de una partícula en un fluido.} Sea $\mathbf{X}_t = (X_t, Y_t, Z_t) \in \mathbb{R}^3$ la posición de una partícula suspendida en un fluido al tiempo $t$ (movimiento browniano). El espacio de estados es $S = \mathbb{R}^3$. Esta es una descripción clásica del movimiento aleatorio de partículas coloidales, estudiado por primera vez por Robert Brown. El espacio parametral puede ser discreto (si observamos la posición en intervalos regulares) o continuo (si consideramos la trayectoria continua de la partícula).
\end{example}

\begin{remark}
La distribución inicial y las probabilidades de transición determinan completamente la evolución probabilística de un proceso estocástico a tiempo discreto. Para un proceso con espacio de estados finito, la matriz de transición es una herramienta algebraica fundamental para estudiar el comportamiento a largo plazo del proceso.
\end{remark}

\subsection*{La ruina del jugador}

Consideremos uno de los problemas más antiguos e importantes de la teoría de procesos estocásticos: \textit{el problema de la ruina del jugador}. Este problema ilustra cómo las probabilidades de transición pueden usarse para resolver preguntas fundamentales sobre la evolución de un proceso.

\textbf{Planteamiento del problema.} Un jugador inicia con una fortuna de $k$ dólares. En cada ronda de juego:
\begin{itemize}
    \item Con probabilidad $p$, el jugador gana \$1 (su fortuna aumenta a $k+1$).
    \item Con probabilidad $q = 1 - p$, el jugador pierde \$1 (su fortuna disminuye a $k-1$).
\end{itemize}

El juego continúa hasta que el jugador se arruina (fortuna = \$0) o alcanza una fortuna objetivo de \$N. Sea $X_n$ la fortuna del jugador después de $n$ rondas, donde $X_0 = k$. El espacio de estados es $S = \{0, 1, 2, \ldots, N\}$, siendo 0 y $N$ \textit{estados absorbentes} (el juego termina cuando se alcanzan).

Las probabilidades de transición de un paso son:
\[
p_{i,i+1} = p, \quad p_{i,i-1} = q, \quad \text{para } i \in \{1, 2, \ldots, N-1\},
\]
y $p_{0,0} = 1$, $p_{N,N} = 1$ (los estados absorbentes permanecen absorbentes).

\textbf{Objetivo.} Calcular $P_k := \mathbb{P}(\text{el jugador alcanza } \$N \mid X_0 = k)$, es decir, la probabilidad de que el jugador gane alcanzando la fortuna objetivo antes de arruinarse.

\textbf{Solución mediante ecuaciones de recurrencia.} 

Razonemos de manera intuitiva. Si el jugador se encuentra en la fortuna $k$ (con $0 < k < N$), en la siguiente ronda ocurre uno de dos eventos:
\begin{enumerate}
    \item Con probabilidad $p$, el jugador gana \$1 y su fortuna pasa a $k+1$.
    \item Con probabilidad $q$, el jugador pierde \$1 y su fortuna pasa a $k-1$.
\end{enumerate}

Por lo tanto, la probabilidad de que el jugador eventualmente alcance \$N partiendo de $k$ depende de lo que suceda en la próxima ronda. Podemos escribir:
\[
P_k = p \cdot P_{k+1} + q \cdot P_{k-1}, \quad \text{para } 0 < k < N.
\]

Con condiciones de frontera:
\[
P_0 = 0 \quad (\text{si el jugador está arruinado, no puede ganar}),
\]
\[
P_N = 1 \quad (\text{si el jugador alcanza la fortuna objetivo, ha ganado}).
\]

Esta es una ecuación de recurrencia lineal homogénea de segundo orden. Su solución depende del valor de $p$.

\textbf{Caso 1: Juego justo} ($p = q = \frac{1}{2}$).

Si $p = \frac{1}{2}$, la ecuación se convierte en:
\[
P_k = \frac{1}{2} P_{k+1} + \frac{1}{2} P_{k-1},
\]
que puede reescribirse como:
\[
P_{k+1} - P_k = P_k - P_{k-1}.
\]

Esto significa que las diferencias consecutivas son constantes. Por lo tanto, $P_k$ es una función lineal de $k$:
\[
P_k = a + bk
\]
para algunas constantes $a$ y $b$. Aplicando las condiciones de frontera:
\[
P_0 = a = 0 \implies a = 0,
\]
\[
P_N = b \cdot N = 1 \implies b = \frac{1}{N}.
\]

Por lo tanto:
\[
\boxed{P_k = \frac{k}{N} \quad \text{(para } p = \tfrac{1}{2}\text{)}.}
\]

Esta es una conclusión importante: \textit{en un juego justo, la probabilidad de ganar es proporcional a la fortuna inicial}.

\textbf{Caso 2: Juego sesgado} ($p \neq \frac{1}{2}$).

Cuando $p \neq q$, buscamos soluciones de la forma $P_k = r^k$. Sustituyendo en la ecuación de recurrencia:
\[
r^k = p \cdot r^{k+1} + q \cdot r^{k-1}.
\]

Dividiendo por $r^{k-1}$:
\[
r = p \cdot r^2 + q,
\]
o equivalentemente:
\[
p r^2 - r + q = 0.
\]

Las soluciones son:
\[
r = \frac{1 \pm \sqrt{1 - 4pq}}{2p} = \frac{1 \pm \sqrt{(1 - 2p)^2}}{2p} = \frac{1 \pm |1 - 2p|}{2p}.
\]

Como $p + q = 1$, tenemos $1 - 4pq = (p - q)^2 = (1 - 2p)^2$, por lo que:
\[
r_1 = 1 \quad \text{y} \quad r_2 = \frac{q}{p}.
\]

La solución general es:
\[
P_k = A + B \left(\frac{q}{p}\right)^k
\]
para constantes $A$ y $B$. Aplicando las condiciones de frontera:
\[
P_0 = A + B = 0 \implies B = -A,
\]
\[
P_N = A + B \left(\frac{q}{p}\right)^N = A \left(1 - \left(\frac{q}{p}\right)^N\right) = 1.
\]

Despejando $A$:
\[
A = \frac{1}{1 - (q/p)^N}.
\]

Por lo tanto:
\[
\boxed{P_k = \frac{1 - (q/p)^k}{1 - (q/p)^N} \quad \text{(para } p \neq \tfrac{1}{2}\text{)}.}
\]

Nótese que esta fórmula recupera el resultado del caso justo cuando $p = \frac{1}{2}$ (tomando el límite adecuado).

\textbf{Interpretación de los resultados.}

\begin{enumerate}
    \item \textbf{Cuando $p > \frac{1}{2}$ (juego favorable).} Tenemos $q/p < 1$, por lo que $(q/p)^k \to 0$ cuando $k \to \infty$. Para un jugador con fortuna inicial $k$ pequeña comparada con $N$, la probabilidad $P_k \approx 1 - (q/p)^k$ es cercana a 1. Esto indica que un jugador con ventaja tiene una alta probabilidad de ganar, especialmente si su objetivo es modesto.
    
    \item \textbf{Cuando $p < \frac{1}{2}$ (juego desfavorable).} Tenemos $q/p > 1$, por lo que $(q/p)^k$ crece exponencialmente con $k$. Incluso si el jugador comienza con una buena fortuna inicial, la probabilidad de ganar disminuye exponencialmente. Para un jugador con $p < \frac{1}{2}$ y en contra de un oponente con una fortuna mucho mayor ($N$ muy grande), la ruina es prácticamente segura.
    
    \item \textbf{El caso justo.} Cuando $p = \frac{1}{2}$, la probabilidad de ganar es simplemente $k/N$, una relación lineal. Un jugador con fortuna $k$ y objetivo $N$ tiene una probabilidad de éxito igual a la proporción de su fortuna inicial respecto al capital total en juego.
\end{enumerate}

\textbf{Observación fundamental.} Este problema demuestra la potencia de analizar procesos estocásticos a través de sus probabilidades de transición. Sin invocar teorías sofisticadas, simplemente planteando ecuaciones de recurrencia basadas en cómo el proceso evoluciona en un paso, hemos obtenido resultados exactos y muy interpretables. Esta es la esencia del enfoque mediante cadenas de Markov: comprender el comportamiento global del sistema a partir del conocimiento local de cómo transiciona entre estados adyacentes.

\section{Cadenas de Markov}

Un \textbf{proceso de Markov} es un proceso estocástico que satisface una propiedad especial de ``falta de memoria'': el futuro del proceso, dado el estado presente, es independiente del pasado. Más formalmente:

\begin{definition}
Decimos que un proceso estocástico $\{X_n : n \in \mathbb{N}\}$ es un \textbf{proceso de Markov} si para todos $n \in \mathbb{N}$ y todos los estados $i_0, i_1, \ldots, i_n, j \in S$:
\[
\mathbb{P}(X_{n+1} = j \mid X_n = i_n, X_{n-1} = i_{n-1}, \ldots, X_0 = i_0) = \mathbb{P}(X_{n+1} = j \mid X_n = i_n),
\]
siempre que las probabilidades condicionales estén bien definidas.
\end{definition}

Una \textbf{cadena de Markov} es un proceso de Markov a tiempo discreto con espacio de estados discreto.

\begin{definition}
Una \textbf{cadena de Markov homogénea} es una cadena de Markov donde las probabilidades de transición no dependen del tiempo $n$, es decir, $\mathbb{P}(X_{n+1} = j \mid X_n = i)$ es la misma para todo $n \geq 0$.
\end{definition}

\begin{remark}
A menos que se indique lo contrario, en este curso supondremos que todas nuestras cadenas de Markov son homogéneas. Esta hipótesis simplifica significativamente el análisis matemático y es apropiada para muchas aplicaciones prácticas.
\end{remark}

\begin{remark}
    La distribución de una cadena de Markov homogénea depende únicamente de la distribución inicial y de las probabilidades de transición, es decir, de la pareja $(\pi,P)$.
\end{remark}

Demostremos que la caminata aleatoria $\{X_n : n \in \mathbb{N}\}$ definida por $X_n = \sum_{i=1}^n Y_i$ (donde los $Y_i$ son i.i.d. con distribución en $S_Y \subseteq \mathbb{R}$) es una cadena de Markov homogénea.

\begin{proposition}
La caminata aleatoria $X_n = \sum_{i=1}^n Y_i$ con $X_0 = 0$ es una cadena de Markov homogénea cuando los $Y_i$ son independientes e idénticamente distribuidas.
\end{proposition}

\begin{proof}
Demostramos dos cosas: primero que es una cadena de Markov, y luego que es homogénea.

\textit{Parte 1: Propiedad de Markov.} 

Sea $n \in \mathbb{N}$ y sean $i_0, i_1, \ldots, i_n, j$ valores específicos en el espacio de estados. Queremos mostrar que:
\[
\mathbb{P}(X_{n+1} = j \mid X_n = i_n, X_{n-1} = i_{n-1}, \ldots, X_0 = i_0) = \mathbb{P}(X_{n+1} = j \mid X_n = i_n).
\]

Por definición, $X_{n+1} = X_n + Y_{n+1}$. Dado que $X_n = i_n$, tenemos:
\[
X_{n+1} = j \iff Y_{n+1} = j - i_n.
\]

Por lo tanto:
\[
\mathbb{P}(X_{n+1} = j \mid X_n = i_n, X_{n-1} = i_{n-1}, \ldots, X_0 = i_0) = \mathbb{P}(Y_{n+1} = j - i_n \mid X_n = i_n, X_{n-1} = i_{n-1}, \ldots, X_0 = i_0).
\]

Ahora bien, la variable aleatoria $Y_{n+1}$ es independiente de todas las variables $Y_1, Y_2, \ldots, Y_n$ por la hipótesis de independencia. Puesto que $X_0, X_1, \ldots, X_n$ son funciones únicamente de $Y_1, Y_2, \ldots, Y_n$, podemos concluir que $Y_{n+1}$ es independiente de $(X_0, X_1, \ldots, X_n)$. En particular:
\[
\mathbb{P}(Y_{n+1} = j - i_n \mid X_n = i_n, X_{n-1} = i_{n-1}, \ldots, X_0 = i_0) = \mathbb{P}(Y_{n+1} = j - i_n).
\]

Asimismo:
\[
\mathbb{P}(X_{n+1} = j \mid X_n = i_n) = \mathbb{P}(Y_{n+1} = j - i_n \mid X_n = i_n) = \mathbb{P}(Y_{n+1} = j - i_n).
\]

Por lo tanto, se cumple la igualdad requerida, y $\{X_n\}$ es una cadena de Markov.

\textit{Parte 2: Homogeneidad.}

Para demostrar que la cadena es homogénea, debemos verificar que la probabilidad de transición $p_{ij}^{[n]} := \mathbb{P}(X_{n+1} = j \mid X_n = i)$ no depende de $n$.

Para cualquier $n$, condicionando en $X_n = i$:
\[
\mathbb{P}(X_{n+1} = j \mid X_n = i) = \mathbb{P}(X_n + Y_{n+1} = j \mid X_n = i) = \mathbb{P}(Y_{n+1} = j - i \mid X_n = i).
\]

Como $Y_{n+1}$ es independiente de $X_n$:
\[
\mathbb{P}(Y_{n+1} = j - i \mid X_n = i) = \mathbb{P}(Y_{n+1} = j - i).
\]

Dado que todos los $Y_i$ tienen la misma distribución, la probabilidad $\mathbb{P}(Y_{n+1} = j - i)$ es idéntica para cualquier $n$. Por lo tanto:
\[
p_{ij} := \mathbb{P}(X_{n+1} = j \mid X_n = i) = \mathbb{P}(Y_1 = j - i)
\]
no depende de $n$, lo que demuestra que la cadena es homogénea. 
\end{proof}

\begin{remark}
Si considerásemos un proceso estocástico $\{Z_n : n \in \mathbb{N}\}$ donde cada $Z_n$ es una variable aleatoria independiente, entonces el proceso sería trivialmente una cadena de Markov homogénea: la probabilidad de transición sería:
\[
\mathbb{P}(Z_{n+1} = j \mid Z_n = i) = \mathbb{P}(Z_{n+1} = j),
\]
que no depende del estado actual $i$. De hecho, para procesos con variables completamente independientes, la propiedad de Markov es automática porque el futuro es independiente de todo el pasado, incluyendo el presente. La especificidad de las cadenas de Markov reside en permitir dependencia entre variables aleatorias sucesivas (a través de las probabilidades de transición) mientras se mantiene la ausencia de memoria más allá del estado presente.
\end{remark}

\subsection*{Ejemplos de cadenas de Markov con matriz de transición explícita}

\begin{example}[12]
\textbf{Cadena de Markov con matriz de transición finita.} Consideremos una cadena de Markov con espacio de estados $S = \{1, 2, 3\}$ y matriz de transición:
\[
P = \begin{pmatrix}
0.5 & 0.3 & 0.2 \\
0.1 & 0.6 & 0.3 \\
0.2 & 0.2 & 0.6
\end{pmatrix},
\]
donde $P_{ij}$ representa la probabilidad $p_{ij} = \mathbb{P}(X_{n+1} = j \mid X_n = i)$. Por ejemplo:
\begin{itemize}
    \item $p_{1,1} = 0.5$: Si el proceso está en estado 1, permanece en estado 1 con probabilidad 0.5.
    \item $p_{1,2} = 0.3$: Si el proceso está en estado 1, transiciona a estado 2 con probabilidad 0.3.
    \item $p_{2,3} = 0.3$: Si el proceso está en estado 2, transiciona a estado 3 con probabilidad 0.3.
\end{itemize}

Cada fila de la matriz suma 1, como corresponde a una matriz estocástica.

Supongamos una distribución inicial $\pi = (0.4, 0.3, 0.3)$, lo que significa:
\[
\mathbb{P}(X_0 = 1) = 0.4, \quad \mathbb{P}(X_0 = 2) = 0.3, \quad \mathbb{P}(X_0 = 3) = 0.3.
\]

A partir de esta distribución inicial y la matriz de transición, podemos calcular la distribución de probabilidades en tiempos posteriores. Por ejemplo, la distribución en tiempo 1 es:
\[
\pi^{(1)} = \pi \cdot P = (0.4, 0.3, 0.3) \cdot P.
\]

Calculando:
\[
\pi_1^{(1)} = 0.4 \cdot 0.5 + 0.3 \cdot 0.1 + 0.3 \cdot 0.2 = 0.2 + 0.03 + 0.06 = 0.29,
\]
\[
\pi_2^{(1)} = 0.4 \cdot 0.3 + 0.3 \cdot 0.6 + 0.3 \cdot 0.2 = 0.12 + 0.18 + 0.06 = 0.36,
\]
\[
\pi_3^{(1)} = 0.4 \cdot 0.2 + 0.3 \cdot 0.3 + 0.3 \cdot 0.6 = 0.08 + 0.09 + 0.18 = 0.35.
\]

Por lo tanto, $\pi^{(1)} = (0.29, 0.36, 0.35)$.
\end{example}

\begin{example}[13]
La \textit{Urna de Ehrenfest} es una cadena de Markov homogénea que representa una descripción matemática simplificada del proceso de difusion de gases o líquidos a través de una membrana. El modelo consiste de dos cajas A y B que contienen un total de $N$ bolas. Seleccionamos al azar una de las $N$ bolas y la colocamos en la otra caja. Sea $X_n$ el número de bolas en la caja $A$ después de $n$ pasos. El espacio de estados es $S = \{0, 1, 2, \ldots, N\}$, y las probabilidades de transición son:
\[
p_{i, i+1} = \frac{N - i}{N} \quad \text{(si seleccionamos una bola de la caja B)},
\]
\[p_{i, i-1} = \frac{i}{N} \quad \text{(si seleccionamos una bola de la caja A)},
\]
y $p_{ij} = 0$ para cualquier otro par $(i, j)$.
\end{example}

\begin{example}[14]
\textbf{La cadena de las horas del reloj.} Sea $X_n$ la hora que marca un reloj después de $n$ horas, donde las horas se representan como $S = \{0, 1, 2, \ldots, 11\}$ (con 0 representando las 12 en punto). El proceso transiciona según la regla determinística:
\[
X_{n+1} = (X_n + 1) \bmod 12.
\]

Esto significa que de cualquier hora $i$, el reloj avanza a la hora $i+1$ (módulo 12) con probabilidad 1. Formalmente, las probabilidades de transición son:
\[
p_{i, (i+1) \bmod 12} = 1 \quad \text{para todo } i \in S,
\]
y $p_{ij} = 0$ para cualquier otro par $(i, j)$.

La matriz de transición es:
\[
P = \begin{pmatrix}
0 & 1 & 0 & 0 & \cdots & 0 \\
0 & 0 & 1 & 0 & \cdots & 0 \\
0 & 0 & 0 & 1 & \cdots & 0 \\
\vdots & \vdots & \vdots & \vdots & \ddots & \vdots \\
1 & 0 & 0 & 0 & \cdots & 0
\end{pmatrix},
\]
donde cada fila contiene exactamente un 1 (en la posición correspondiente a la siguiente hora) y ceros en todas partes.

Este es un ejemplo de cadena de Markov con estructura \textit{cíclica}. Si el reloj comienza en las 3 en punto ($X_0 = 3$), después de 1 hora marca las 4 ($X_1 = 4$), después de 2 horas marca las 5 ($X_2 = 5$), y después de 12 horas regresa a las 3 ($X_{12} = 3$). De manera general, $X_n = (X_0 + n) \bmod 12$.

Aunque el proceso es determinístico (no hay aleatoriedad en las transiciones), sigue siendo una cadena de Markov válida: el futuro depende únicamente del estado presente, no del pasado.
\end{example}

Para aumentar nuestro entendimiento de la propiedad de Markov, probaremos la siguiente proposición, que auda en la interpretación del futuro independiente del pasado dado el presente.

\begin{proposition}
    Sea $X$ una cadena de Markov homogénea $(\pi,P)$, con espacio de estados $S$, $i,j,k\in S$  y $0\leq m \leq n-1$. Se tiene que 
    $$\mathbb{P}(X_{n+1}=j | X_n = i, X_{m} = k) = \mathbb{P}(X_{n+1}=j | X_n = i).$$ 

\end{proposition}

\section[Probabilidades de Transición en n Pasos]{Probabilidades de Transición en $n$ Pasos}

Hasta ahora hemos estudiado las probabilidades de transición de un paso, definidas como $p_{ij} = \mathbb{P}(X_{n+1} = j \mid X_n = i)$. Sin embargo, en muchas aplicaciones es fundamental conocer la probabilidad de que la cadena pase del estado $i$ al estado $j$ transcurridos exactamente $n$ pasos, independientemente del instante inicial.

\begin{definition}
Sea $\{X_n : n \in \mathbb{N}\}$ una cadena de Markov homogénea con espacio de estados $S$. Definimos la \textbf{probabilidad de transición en $n$ pasos} como:
\[
p_{ij}^{(n)} = \mathbb{P}(X_{m+n} = j \mid X_m = i),
\]
para cualesquiera $i, j \in S$ y $m, n \in \mathbb{N}$. La homogeneidad de la cadena garantiza que esta probabilidad es independiente del instante inicial $m$.
\end{definition}

\begin{remark}
Por supuesto, para $n = 1$, recuperamos las probabilidades de transición de un paso:
\[
p_{ij}^{(1)} = p_{ij}.
\]
Asimismo, para $n = 0$, se define:
\[
p_{ij}^{(0)} = \begin{cases} 1 & \text{si } i = j, \\ 0 & \text{si } i \neq j. \end{cases}
\]
Esta es la matriz identidad, lo que tiene sentido: sin dar ningún paso, la cadena permanece en su estado actual con certeza.
\end{remark}

Cuando el espacio de estados es finito, estas probabilidades pueden organizarse en una \textbf{matriz de transición en $n$ pasos}:
\[
P^{(n)} = (p_{ij}^{(n)})_{i,j \in S},
\]
donde $P^{(1)} = P$ es la matriz de transición estándar y $P^{(0)} = I$ es la matriz identidad.

\begin{example}[15]
Consideremos la cadena de Markov del Ejemplo 12 con matriz de transición:
\[
P = \begin{pmatrix}
0.5 & 0.3 & 0.2 \\
0.1 & 0.6 & 0.3 \\
0.2 & 0.2 & 0.6
\end{pmatrix}.
\]

¿Cuál es la probabilidad de que la cadena transite del estado 1 al estado 2 en exactamente 2 pasos? Esta es la cantidad $p_{1,2}^{(2)}$. Para calcularla, podemos identificar todos los caminos de longitud 2 que comienzan en el estado 1 y terminan en el estado 2:
\begin{itemize}
    \item Camino $1 \to 1 \to 2$: Probabilidad $p_{1,1} \cdot p_{1,2} = 0.5 \times 0.3 = 0.15$.
    \item Camino $1 \to 2 \to 2$: Probabilidad $p_{1,2} \cdot p_{2,2} = 0.3 \times 0.6 = 0.18$.
    \item Camino $1 \to 3 \to 2$: Probabilidad $p_{1,3} \cdot p_{3,2} = 0.2 \times 0.2 = 0.04$.
\end{itemize}

Por lo tanto:
\[
p_{1,2}^{(2)} = 0.15 + 0.18 + 0.04 = 0.37.
\]

Alternativamente, podemos calcular $P^{(2)} = P \cdot P = P^2$:
\[
P^2 = \begin{pmatrix}
0.5 & 0.3 & 0.2 \\
0.1 & 0.6 & 0.3 \\
0.2 & 0.2 & 0.6
\end{pmatrix} \begin{pmatrix}
0.5 & 0.3 & 0.2 \\
0.1 & 0.6 & 0.3 \\
0.2 & 0.2 & 0.6
\end{pmatrix}.
\]

El elemento $(1,2)$ de $P^2$ es:
\[
(P^2)_{1,2} = 0.5 \times 0.3 + 0.3 \times 0.6 + 0.2 \times 0.2 = 0.15 + 0.18 + 0.04 = 0.37,
\]
que coincide con nuestro cálculo anterior.
\end{example}

\subsection*{La Ecuación de Chapman-Kolmogorov}

La ecuación de Chapman-Kolmogorov es una relación fundamental que conecta probabilidades de transición en diferentes números de pasos. Esta ecuación expresa una intuición natural: para ir del estado $i$ al estado $j$ en $n + m$ pasos, necesariamente pasamos por algún estado intermedio $k$ después de los primeros $n$ pasos.

\begin{theorem}[Ecuación de Chapman-Kolmogorov]
Sea $\{X_n : n \in \mathbb{N}\}$ una cadena de Markov homogénea con espacio de estados $S$. Entonces, para cualesquiera enteros $n, m \geq 0$ y estados $i, j \in S$:
\[
p_{ij}^{(n+m)} = \sum_{k \in S} p_{ik}^{(n)} \cdot p_{kj}^{(m)}.
\]

En notación matricial:
\[
P^{(n+m)} = P^{(n)} \cdot P^{(m)},
\]
donde el producto es el producto usual de matrices.
\end{theorem}

\begin{proof}
Demostraremos la ecuación utilizando la ley de probabilidad total y la propiedad de Markov.

Sean $n, m \geq 0$ y consideremos los estados $i, j \in S$. Por definición:
\[
p_{ij}^{(n+m)} = \mathbb{P}(X_{n+m} = j \mid X_0 = i).
\]

Condicionamos sobre el estado de la cadena después de $n$ pasos, $X_n$. Por la ley de probabilidad total:
\[
\mathbb{P}(X_{n+m} = j \mid X_0 = i) = \sum_{k \in S} \mathbb{P}(X_{n+m} = j \mid X_n = k, X_0 = i) \cdot \mathbb{P}(X_n = k \mid X_0 = i).
\]

Ahora aplicamos la propiedad de Markov (Definición de cadena de Markov). Dado que la cadena es Markoviana y el estamos en el estado $X_n = k$, el futuro (desde el tiempo $n$ en adelante) es independiente del pasado (antes del tiempo $n$). Específicamente:
\[
\mathbb{P}(X_{n+m} = j \mid X_n = k, X_0 = i) = \mathbb{P}(X_{n+m} = j \mid X_n = k).
\]

Por lo tanto:
\[
p_{ij}^{(n+m)} = \sum_{k \in S} \mathbb{P}(X_{n+m} = j \mid X_n = k) \cdot \mathbb{P}(X_n = k \mid X_0 = i).
\]

Reconocemos en esta expresión:
\[
\mathbb{P}(X_{n+m} = j \mid X_n = k) = \mathbb{P}(X_{m} = j \mid X_0 = k) = p_{kj}^{(m)},
\]
y
\[
\mathbb{P}(X_n = k \mid X_0 = i) = p_{ik}^{(n)},
\]
por la homogeneidad de la cadena. Por lo tanto:
\[
p_{ij}^{(n+m)} = \sum_{k \in S} p_{ik}^{(n)} \cdot p_{kj}^{(m)},
\]
que es exactamente la ecuación de Chapman-Kolmogorov.

En forma matricial, esta ecuación correspond a la multiplicación de matrices, es decir, $P^{(n+m)} = P^{(n)} \cdot P^{(m)}$.
En particular, la siguiente desigualdad será utilizada más adelante: para cualesquiera estados $i, j, k$ y para $0 \leq r \leq n$, se tiene que:
\[
p_{ij}^{(n)} \geq  p_{ik}^{(r)} p_{kj}^{(n-r)}.
\]

\end{proof}

\begin{remark}
La ecuación de Chapman-Kolmogorov es fundamental porque implica que:
\[
P^{(n)} = \underbrace{P \times P \times \cdots \times P}_{n \text{ veces}} = P^n,
\]
es decir, la matriz de transición en $n$ pasos es simplemente la $n$-ésima potencia de la matriz de transición de un paso $P$. Esta es una consecuencia inmediata de la ecuación de Chapman-Kolmogorov aplicada recursivamente:
\[
P^{(n)} = P^{(n-1+1)} = P^{(n-1)} \cdot P^{(1)} = P^{(n-1)} \cdot P.
\]

Por inducción, obtenemos $P^{(n)} = P^n$.
\end{remark}

\begin{example}[16]
Retomemos la cadena de Markov del Ejemplo 12. Queremos calcular $p_{1,2}^{(3)}$, la probabilidad de ir del estado 1 al estado 2 en exactamente 3 pasos.

Por la ecuación de Chapman-Kolmogorov con $n = 2$ y $m = 1$:
\[
p_{1,2}^{(3)} = \sum_{k=1}^{3} p_{1,k}^{(2)} \cdot p_{k,2}^{(1)}.
\]

Primero necesitamos $P^{(2)} = P^2$. Ya calculamos este en el Ejemplo 15:
\[
P^2 = \begin{pmatrix}
0.29 + 0.11 & 0.36 & 0.35 \\
\cdot & \cdot & \cdot \\
\cdot & \cdot & \cdot
\end{pmatrix}
\quad \text{(mostramos solo la fila relevante con detalle)}.
\]

Calculemos $P^2$ completamente. El elemento $(i,j)$ de $P^2$ es $\sum_k p_{ik} p_{kj}$. Después del cálculo:
\begin{align*}
P^2 &= \begin{pmatrix}
0.5 \cdot 0.5 + 0.3 \cdot 0.1 + 0.2 \cdot 0.2 & 0.5 \cdot 0.3 + 0.3 \cdot 0.6 + 0.2 \cdot 0.2 & 0.5 \cdot 0.2 + 0.3 \cdot 0.3 + 0.2 \cdot 0.6 \\
0.1 \cdot 0.5 + 0.6 \cdot 0.1 + 0.3 \cdot 0.2 & 0.1 \cdot 0.3 + 0.6 \cdot 0.6 + 0.3 \cdot 0.2 & 0.1 \cdot 0.2 + 0.6 \cdot 0.3 + 0.3 \cdot 0.6 \\
0.2 \cdot 0.5 + 0.2 \cdot 0.1 + 0.6 \cdot 0.2 & 0.2 \cdot 0.3 + 0.2 \cdot 0.6 + 0.6 \cdot 0.2 & 0.2 \cdot 0.2 + 0.2 \cdot 0.3 + 0.6 \cdot 0.6
\end{pmatrix}\\
&= \begin{pmatrix}
0.25 + 0.03 + 0.04 & 0.15 + 0.18 + 0.04 & 0.10 + 0.09 + 0.12 \\
0.05 + 0.06 + 0.06 & 0.03 + 0.36 + 0.06 & 0.02 + 0.18 + 0.18 \\
0.10 + 0.02 + 0.12 & 0.06 + 0.12 + 0.12 & 0.04 + 0.06 + 0.36
\end{pmatrix}\\
&= \begin{pmatrix}
0.32 & 0.37 & 0.31 \\
0.17 & 0.45 & 0.38 \\
0.24 & 0.30 & 0.46
\end{pmatrix}.
\end{align*}

Ahora aplicamos Chapman-Kolmogorov:
\[
p_{1,2}^{(3)} = p_{1,1}^{(2)} \cdot p_{1,2} + p_{1,2}^{(2)} \cdot p_{2,2} + p_{1,3}^{(2)} \cdot p_{3,2}
\]
\[
= 0.32 \times 0.3 + 0.37 \times 0.6 + 0.31 \times 0.2 = 0.096 + 0.222 + 0.062 = 0.380.
\]

Alternativamente, podemos calcular $P^3 = P^2 \times P$ y leer el elemento $(1,2)$.
\end{example}

\begin{remark}
La ecuación de Chapman-Kolmogorov no es únicamente un resultado teórico interesante; tiene profundas implicaciones prácticas. Al expresar $P^{(n)} = P^n$, nos proporciona un algoritmo eficiente (usando exponenciación rápida de matrices) para calcular probabilidades de transición en un gran número de pasos sin necesidad de enumerar todos los caminos posibles. Esto es especialmente relevante en aplicaciones computacionales donde frecuentemente nos interesa el comportamiento a largo plazo de cadenas de Markov.
\end{remark}

\subsection*{Diagonalización de la Matriz de Transición}

Para calcular $P^n$ de manera eficiente, especialmente cuando $n$ es grande, podemos utilizar la diagonalización de la matriz $P$. El siguiente teorema es fundamental en el álgebra lineal y tiene aplicaciones directas al estudio de cadenas de Markov.

\begin{theorem}
Sea $P$ una matriz cuadrada $m \times m$ con espectro (conjunto de valores propios) $\Lambda = \{\lambda_1, \lambda_2, \ldots, \lambda_m\}$, donde $\lambda_i$ son valores propios (no necesariamente distintos) y sean $\mathbf{v}_1, \mathbf{v}_2, \ldots, \mathbf{v}_m$ los correspondientes vectores propios por la derecha.\footnote{Un vector propio derecho satisface $P \mathbf{v}_i = \lambda_i \mathbf{v}_i$.}

Si $P$ es diagonalizable\footnote{Una condición suficiente para que una matriz sea diagonalizable es que todos sus valores propios sean distintos.}, entonces existe una matriz $V$ cuyas columnas son los vectores propios $\mathbf{v}_1, \mathbf{v}_2, \ldots, \mathbf{v}_m$, y una matriz $D = \text{diag}(\lambda_1, \lambda_2, \ldots, \lambda_m)$ diagonal de valores propios, tales que:
\[
P = V D V^{-1}.
\]

En consecuencia, para el cálculo de potencias:
\[
P^n = V D^n V^{-1},
\]
donde $D^n = \text{diag}(\lambda_1^n, \lambda_2^n, \ldots, \lambda_m^n)$ es trivial de calcular.
\end{theorem}

\begin{remark}
Para una matriz cuadrada $P$ (en particular, una matriz de transición), existen tres propiedades fundamentales relacionadas con los valores propios:

\begin{enumerate}
    \item Si la matriz es estocástica, tiene al menos un valor propio igual a 1, lo cual está relacionado con la existencia de distribuciones estacionaria (por verse en las secciones siguientes). 

    \item La \textbf{traza} de $P$ (suma de los elementos diagonales) es igual a la suma de los valores propios:
    \[
    \text{tr}(P) = \sum_{i=1}^m p_{ii} = \sum_{i=1}^m \lambda_i.
    \]
    
    \item El \textbf{determinante} de $P$ es el producto de los valores propios:
    \[
    \det(P) = \prod_{i=1}^m \lambda_i.
    \]
\end{enumerate}
En los casos $2 \times 2$ o $3 \times 3$, es posible obtener los valores propios resolviendo el sistema de ecuaciones que resulta de estas tres propiedades, sin necesidad de calcular explícitamente el polinomio característico.
\end{remark}

\begin{example}[17]
\textbf{Cadena de Markov con dos estados: caso general.} Consideremos una cadena de Markov con espacio de estados $S = \{1, 2\}$ y matriz de transición:
\[
P = \begin{pmatrix}
1 - \alpha & \alpha \\
\beta & 1 - \beta
\end{pmatrix},
\]
donde $0 < \alpha, \beta < 1$. Esta es una cadena arbitraria simétrica o asimétrica, dependiendo de los valores de $\alpha$ y $\beta$.

\textbf{Paso 1: Calcular los valores propios.}

Para esto podemos encontrar las raíces del polinomio característico $\det(P - \lambda I) = 0$ o hacer caso a las observaciones anteriores. Al ser el determinante es $\det(P) = (1 - \alpha)(1 - \beta) - \alpha \beta = 1 - \alpha - \beta$ y 1 un valor propio, el otro tiene que necesariamente ser $1 - \alpha - \beta$.

\textbf{Paso 2: Calcular los vectores propios.}

Para $\lambda_1 = 1$:
\[
\begin{pmatrix}
1 - \alpha - 1 & \alpha \\
\beta & 1 - \beta - 1
\end{pmatrix} \begin{pmatrix} v_1 \\ v_2 \end{pmatrix} = \begin{pmatrix} 0 \\ 0 \end{pmatrix}
\]
\[
\begin{pmatrix}
-\alpha & \alpha \\
\beta & -\beta
\end{pmatrix} \begin{pmatrix} v_1 \\ v_2 \end{pmatrix} = \begin{pmatrix} 0 \\ 0 \end{pmatrix}.
\]

De la primera ecuación: $-\alpha v_1 + \alpha v_2 = 0 \implies v_1 = v_2$. Un vector propio es $\mathbf{v}_1 = \begin{pmatrix} 1 \\ 1 \end{pmatrix}$.

Para $\lambda_2 = 1 - \alpha - \beta$:
\[
\begin{pmatrix}
1 - \alpha - (1 - \alpha - \beta) & \alpha \\
\beta & 1 - \beta - (1 - \alpha - \beta)
\end{pmatrix} \begin{pmatrix} v_1 \\ v_2 \end{pmatrix} = \begin{pmatrix} 0 \\ 0 \end{pmatrix}
\]
\[
\begin{pmatrix}
\beta & \alpha \\
\beta & \alpha
\end{pmatrix} \begin{pmatrix} v_1 \\ v_2 \end{pmatrix} = \begin{pmatrix} 0 \\ 0 \end{pmatrix}.
\]

De la primera ecuación: $\beta v_1 + \alpha v_2 = 0 \implies v_2 = -\frac{\beta}{\alpha} v_1$. Un vector propio es $\mathbf{v}_2 = \begin{pmatrix} 1 \\ -\beta/\alpha \end{pmatrix}$ o, reescalando, $\mathbf{v}_2 = \begin{pmatrix} \alpha \\ -\beta \end{pmatrix}$.

\textbf{Paso 3: Formar la matriz $V$ y su inversa.}

\[
V = \begin{pmatrix} 1 & \alpha \\ 1 & -\beta \end{pmatrix}.
\]

El determinante es $\det(V) = 1 \cdot (-\beta) - \alpha \cdot 1 = -\beta - \alpha = -(\alpha + \beta)$.

\[
V^{-1} = \frac{1}{-(\alpha + \beta)} \begin{pmatrix} -\beta & -\alpha \\ -1 & 1 \end{pmatrix} = \frac{1}{\alpha + \beta} \begin{pmatrix} \beta & \alpha \\ 1 & -1 \end{pmatrix}.
\]

\textbf{Paso 4: Expresión para $P^n$.}

\[
P^n = V D^n V^{-1} = \begin{pmatrix} 1 & \alpha \\ 1 & -\beta \end{pmatrix} \begin{pmatrix} 1 & 0 \\ 0 & (1-\alpha-\beta)^n \end{pmatrix} \frac{1}{\alpha + \beta} \begin{pmatrix} \beta & \alpha \\ 1 & -1 \end{pmatrix}.
\]

Calculando:
\[
V D^n = \begin{pmatrix} 1 & \alpha(1-\alpha-\beta)^n \\ 1 & -\beta(1-\alpha-\beta)^n \end{pmatrix}.
\]

Multiplicando por $V^{-1}$:
\[
P^n = \frac{1}{\alpha + \beta} \begin{pmatrix} 1 & \alpha(1-\alpha-\beta)^n \\ 1 & -\beta(1-\alpha-\beta)^n \end{pmatrix} \begin{pmatrix} \beta & \alpha \\ 1 & -1 \end{pmatrix}.
\]

El elemento $(1,1)$ de $P^n$ es:
\[
(P^n)_{1,1} = \frac{1}{\alpha + \beta} \left[ \beta + \alpha(1-\alpha-\beta)^n \right].
\]

De manera similar, podemos calcular los demás elementos:
\[
\boxed{P^n = \frac{1}{\alpha + \beta} \begin{pmatrix} 
\beta + \alpha(1-\alpha-\beta)^n & \alpha - \alpha(1-\alpha-\beta)^n \\
\beta - \beta(1-\alpha-\beta)^n & \alpha + \beta(1-\alpha-\beta)^n
\end{pmatrix}.}
\]

\textbf{Comportamiento asintótico cuando $n \to \infty$:}

Dado que $0 < \alpha, \beta < 1$, tenemos $0 < \alpha + \beta < 2$, y por lo tanto $1 - \alpha - \beta \in (-1, 1)$. Esto implica que $(1 - \alpha - \beta)^n \to 0$ cuando $n \to \infty$.

Por consiguiente:
\[
\lim_{n \to \infty} P^n = \frac{1}{\alpha + \beta} \begin{pmatrix} 
\beta & \alpha \\
\beta & \alpha
\end{pmatrix}.
\]

Este es un fenómeno notablemente importante: \textit{independientemente del estado inicial}, después de un número muy grande de pasos, la probabilidad de encontrarse en el estado $j$ es esencialmente la misma. Esto significa que la ``memoria'' del estado inicial se disipa con el tiempo, y la cadena alcanza un régimen estacionario donde la distribución de probabilidades se estabiliza.

Esto establece que la distribución límite satisface una ecuación de equilibrio: las probabilidades no cambian de un paso al siguiente. Este comportamiento es intuitivo desde una perspectiva física: después de suficiente tiempo, el sistema ``olvida'' de dónde comenzó y se comporta como si estuviera en un estado de equilibrio.
En las secciones siguientes especificaremos las condiciones bajo las cuales una cadena de Markov garantiza la existencia y unicidad de una distribución estacionaria, así como condiciones que aseguran la convergencia.
\end{example}



\subsection*{Cálculo de la inversa de una matriz}

Para completar el cálculo de $P^n$ mediante diagonalización, frecuentemente necesitamos calcular la inversa de la matriz $V$ de vectores propios. A continuación recordamos el método general usando cofactores y determinantes.

\begin{definition}
Sea $A$ una matriz cuadrada $m \times m$. La \textbf{matriz de cofactores} $C$ se define como:
\[
C_{ij} = (-1)^{i+j} M_{ij},
\]
donde $M_{ij}$ es el determinante del menor obtenido al eliminar la fila $i$ y la columna $j$ de $A$. La \textbf{matriz adjunta} (o \textbf{matriz de adjuntos}) es la transpuesta de la matriz de cofactores:
\[
\text{adj}(A) = C^T.
\]
\end{definition}

\begin{theorem}
Sea $A$ una matriz cuadrada con $\det(A) \neq 0$. Entonces:
\[
A^{-1} = \frac{1}{\det(A)} \text{adj}(A).
\]
\end{theorem}

\begin{remark}
Este método es especialmente útil para matrices pequeñas (por ejemplo $3 \times 3$). Para matrices grandes, métodos computacionales como eliminación gaussiana o descomposición LU suelen ser más eficientes.
\end{remark}

\textbf{Caso particular: Matriz $3 \times 3$.}

Para una matriz $A = \begin{pmatrix} a_{11} & a_{12} & a_{13} \\ a_{21} & a_{22} & a_{23} \\ a_{31} & a_{32} & a_{33} \end{pmatrix}$, primero calculamos:
\[
\det(A) = a_{11}(a_{22}a_{33} - a_{23}a_{32}) - a_{12}(a_{21}a_{33} - a_{23}a_{31}) + a_{13}(a_{21}a_{32} - a_{22}a_{31}).
\]

Luego, la matriz de cofactores es:
\[
C = \begin{pmatrix}
+(a_{22}a_{33} - a_{23}a_{32}) & -(a_{21}a_{33} - a_{23}a_{31}) & +(a_{21}a_{32} - a_{22}a_{31}) \\
-(a_{12}a_{33} - a_{13}a_{32}) & +(a_{11}a_{33} - a_{13}a_{31}) & -(a_{11}a_{32} - a_{12}a_{31}) \\
+(a_{12}a_{23} - a_{13}a_{22}) & -(a_{11}a_{23} - a_{13}a_{21}) & +(a_{11}a_{22} - a_{12}a_{21})
\end{pmatrix}.
\]

La matriz adjunta es $\text{adj}(A) = C^T$, y finalmente:
\[
A^{-1} = \frac{1}{\det(A)} \text{adj}(A).
\]

\begin{example}[18]
\textbf{Cadena de Markov con tres estados.} Consideremos una cadena de Markov con espacio de estados $S = \{1, 2, 3\}$ y matriz de transición:
\[
P = \begin{pmatrix}
0 & \tfrac{1}{2} & \tfrac{1}{2} \\
\tfrac{1}{4} & \tfrac{1}{2} & \tfrac{1}{4} \\
\tfrac{1}{4} & \tfrac{1}{4} & \tfrac{1}{2}
\end{pmatrix}.
\]
\end{example}

\section{Comunicación entre Estados}

Una pregunta fundamental en la teoría de cadenas de Markov es: ¿puede la cadena alcanzar un estado partiendo desde otro? Esta pregunta nos conduce a las nociones de \textbf{accesibilidad} y \textbf{comunicación}, conceptos que clasifican los estados de acuerdo con su ``conectividad'' en el grafo de transiciones.

\begin{definition}
Sea $\{X_n : n \in \mathbb{N}\}$ una cadena de Markov homogénea con espacio de estados $S$. Decimos que un estado $j$ es \textbf{accesible} desde un estado $i$ (denotado $i \to j$) si existe un entero $n \geq 0$ tal que:
\[
\mathbb{P}(X_n = j \mid X_0 = i) = p_{ij}^{(n)} > 0.
\]

Equivalentemente, existe un camino (posiblemente de longitud cero) que conecta $i$ con $j$ con probabilidad positiva.
\end{definition}

\begin{definition}
Dos estados $i$ y $j$ \textbf{se comunican} (denotado $i \leftrightarrow j$) si son mutuamente accesibles, es decir, si $i \to j$ y $j \to i$.
\end{definition}

\begin{remark}
La relación ``se comunica con'' es una relación de equivalencia: es reflexiva ($i \leftrightarrow i$ para todo $i$), simétrica ($i \leftrightarrow j \implies j \leftrightarrow i$), y transitiva ($i \leftrightarrow j$ y $j \leftrightarrow k$ implica $i \leftrightarrow k$). 

Por lo tanto, la relación de comunicación particiona el espacio de estados $S$ en \textbf{clases de comunicación} disjuntas. Dos estados pertenecen a la misma clase si y sólo si se comunican entre sí.
\end{remark}

\begin{example}[19]
\textbf{Movilidad social revisitada.} Reconsideremos la cadena de Markov del Ejemplo 4 que modela la movilidad socioeconómica con estados $S = \{B, M, A\}$ (baja, media, alta). Supongamos que la matriz de transición es:
\[
P = \begin{pmatrix}
0.6 & 0.3 & 0.1 \\
0.2 & 0.5 & 0.3 \\
0 & 0 & 1
\end{pmatrix},
\]
donde las filas y columnas corresponden a los estados $B$, $M$ y $A$, respectivamente.

Analicemos la accesibilidad:
\begin{itemize}
    \item Desde $B$ podemos alcanzar $M$ con probabilidad $p_{BM} = 0.3 > 0$, de modo que $B \to M$.
    \item Desde $M$ podemos alcanzar $A$ con probabilidad $p_{MA} = 0.3 > 0$, de modo que $M \to A$.
    \item Desde $A$ no podemos alcanzar $B$ o $M$ 
\end{itemize}

En consecuencia, tenemos dos clases de comunicación:
\begin{enumerate}
    \item Clase 1: $\{B, M\}$ donde $B \leftrightarrow M$ (ambos se comunican).
    \item Clase 2: $\{A\}$ que es una clase unipuntual (estado absorbente).
\end{enumerate}

Observe que una vez que la cadena entra la tercera clase $\{A\}$ (estado de clase alta), permanece ahí indefinidamente con probabilidad 1.
\end{example}

\begin{definition}
Una cadena de Markov se llama \textbf{irreducible} si todos sus estados pertenecen a una única clase de comunicación, es decir, si cualesquiera dos estados $i, j \in S$ se comunican mutuamente ($i \leftrightarrow j$). De forma equivalente, todo estado es accesible desde cualquier otro estado.
\end{definition}

\begin{example}[20]
\textbf{Cadena con varias clases de comunicación.} Consideremos una cadena de Markov con espacio de estados $S = \{1, 2, 3, 4, 5, 6\}$ y matriz de transición:
\[
P = \begin{pmatrix}
0.5 & 0.5 & 0 & 0 & 0 & 0 \\
0.3 & 0.7 & 0 & 0 & 0 & 0 \\
0 & 0 & 1 & 0 & 0 & 0 \\
0 & 0 & 0.4 & 0 & 0.6 & 0 \\
0 & 0 & 0 & 0 & 1 & 0 \\
0 & 0 & 0 & 0 & 0 & 1
\end{pmatrix}.
\]

Analicemos la estructura de comunicación:
\begin{itemize}
    \item \textbf{Clase 1}: Estados $\{1, 2\}$. Desde el estado 1 podemos alcanzar el estado 2 con $p_{1,2} = 0.5 > 0$, y desde el estado 2 podemos alcanzar el estado 1 con $p_{2,1} = 0.3 > 0$. Por lo tanto, $1 \leftrightarrow 2$. Observe que desde estos estados no podemos escapar: $p_{1,3} = p_{1,4} = p_{1,5} = p_{1,6} = 0$ y $p_{2,3} = p_{2,4} = \cdots = p_{2,6} = 0$.
    
    \item \textbf{Clase 2}: Estado $\{3\}$. Este es un estado absorbente: $p_{3,3} = 1$ y $p_{3,j} = 0$ para $j \neq 3$. Una vez que la cadena entra en este estado, permanece ahí para siempre. 
    
    \item \textbf{Clase 3}: Estado $\{4\}$. Desde el estado 4 podemos alcanzar el estado 5 con probabilidad $p_{4,5} = 0.6 > 0$, pero desde el estado 5 tenemos $p_{5,5} = 1$ y $p_{5,4} = 0$, lo que significa que una vez en el estado 5, la cadena permanece allí. Así pues, $4 \to 5$ pero NO $5 \to 4$, de modo que 4 y 5 NO se comunican. Como $\{4\}$ no se comunica con otros estados, forma entonces su propia clase de comunicación.
    \item \textbf{Clase 4}: Estado $\{5\}$. Este es un estado absorbente: $p_{5,5} = 1$ y $p_{5,j} = 0$ para $j \neq 5$. 
    \item \textbf{Clase 5}: Estado $\{6\}$. Este es un estado absorbente: $p_{6,6} = 1$.
\end{itemize}


Esta cadena es \textbf{reducible} porque posee cinco clases de comunicación.
\end{example}

\begin{example}[21]
\textbf{Caminata aleatoria simétrica.} La caminata aleatoria simétrica del Ejemplo 2 con espacio de estados $S = \mathbb{Z}$ es un ejemplo de cadena irreducible. Desde cualquier entero $i$, podemos alcanzar cualquier otro entero $j$ seguindo una secuencia de pasos hacia arriba (incrementos de $+1$) y hacia abajo (decrementos de $-1$). Por ejemplo, desde el estado 0 podemos alcanzar el estado 5 mediante la secuencia $0 \to 1 \to 2 \to 3 \to 4 \to 5$, cada transición ocurriendo con probabilidad $\frac{1}{2}$. De hecho, existen múltiples caminos desde 0 hasta 5, por lo que $0 \to 5$ (y mutuamente $5 \to 0$).

Esta irreducibilidad es una propiedad estructural fundamental de la caminata aleatoria simétrica: la cadena ``explora'' potencialmente todo el espacio de estados.
\end{example}

\begin{example}[22]
\textbf{Cadena de las horas del reloj.} Recordemos el Ejemplo 14 de la cadena de las horas. Desde cualquier hora $i$, podemos alcanzar cualquier otra hora $j$ avanzando $j - i \pmod{12}$ horas. Por ejemplo, desde las 3 podemos alcanzar las 7 avanzando 4 horas ($3 \to 4 \to 5 \to 6 \to 7$), o avanzando 16 horas ($3 \to 4 \to \cdots \to 19 \equiv 7 \pmod{12}$). Además, podemos regresar desde las 7 a las 3 avanzando 8 horas. Por lo tanto, dos estados cualesquiera se comunican, y la cadena es irreducible.
\end{example}

\section{Distribuciones Periódicas: Período de un Estado}

Hemos visto que ciertos estados pueden ser accesibles desde otros, pero existe una propiedad adicional que afecta la estructura de la cadena: el \textbf{período} de un estado. El período mide la periodicidad con la cual la cadena puede regresar a un estado dado.

\begin{definition}
El \emph{periodo} de un estado $i$ es un número entero no negativo,
denotado por $d_i$, y definido como
\[
d_i = \gcd \{\, n \ge 1 : p_{ii}^{(n)} > 0 \,\},
\]
donde $\gcd$ denota el máximo común divisor. Si $p_{ii}^{(n)} = 0$ para toda $n \ge 1$, se define $d_i = 0$. En particular, se dice que un estado $i$ es \emph{aperiódico} si $d_i = 1$.
Cuando $d_i = k \ge 2$, se dice que $i$ es \emph{periódico de periodo $k$}.
\end{definition}



\begin{example}[22]
\textbf{Cadena de las horas cíclica.} Recordemos la cadena del Ejemplo 14, donde $X_{n+1} = (X_n + 1) \bmod 12$. Debido a que cada estado se alcanza incrementando el reloj en 1 unidad, tenemos:
\begin{itemize}
    \item La cadena retorna al estado 0 (las 12 en punto) después de exactamente 12 pasos: $p_{0,0}^{(12)} = 1$.
    \item La cadena no retorna al estado 0 en ningún otro número de pasos: $p_{0,0}^{(n)} = 0$ para $n \neq 12, 24, 36, \ldots$
    \item Por lo tanto, el conjunto de retorno es $\{12, 24, 36, \ldots\} = \{12k : k \geq 1\}$.
    \item El período es $d_0 = \gcd(12, 24, 36, \ldots) = 12$.
\end{itemize}

De forma similar, para cualquier hora $i$, el período es $d_i = 12$. La cadena es \textbf{periódica con período 12}.
\end{example}

\begin{example}[23]
\textbf{Caminata aleatoria simétrica.} Para la caminata aleatoria simétrica del Ejemplo 2, considere un estado particular, digamos el 0. Tenemos:
\begin{itemize}
    \item En 1 paso: $p_{0,0}^{(1)} = 0$ (no podemos regresar en 1 paso puesto que cada paso es $\pm 1$).
    \item En 2 pasos: $p_{0,0}^{(2)} = \mathbb{P}(Y_1 = 1, Y_2 = -1) + \mathbb{P}(Y_1 = -1, Y_2 = 1) = \frac{1}{4} + \frac{1}{4} = \frac{1}{2} > 0$.
    \item En 3 pasos: $p_{0,0}^{(3)} = 0$ (imposible regresar en un número impar de pasos).
    \item En 4 pasos: $p_{0,0}^{(4)} > 0$ (retornos como $+1, +1, -1, -1$ o $+1, -1, +1, -1$, etc.).
\end{itemize}

El conjunto de retorno es $\{2, 4, 6, 8, \ldots\}$ (todos los números pares), y el período es:
\[
d_0 = \gcd(2, 4, 6, 8, \ldots) = 2.
\]
\end{example}

\begin{proposition}
El período es una \textbf{propiedad de clase}: si dos estados $i$ y $j$ se comunican ($i \leftrightarrow j$), entonces tienen el mismo período: $d_i = d_j$.
\end{proposition}

\begin{proof}
Supongamos que $i \leftrightarrow j$. Entonces existen $m,n \ge 1$ tales que
\[
p_{ij}^{(m)} > 0,
\qquad
p_{ji}^{(n)} > 0.
\]

Definimos
\[
A_i=\{r\ge1:\; p_{ii}^{(r)}>0\},
\qquad
A_j=\{r\ge1:\; p_{jj}^{(r)}>0\},
\]
y $d_i=\gcd(A_i)$, $d_j=\gcd(A_j)$.

\medskip

\textbf{Paso 1: $d_j \mid d_i$.}

Sea $r\in A_i$. Entonces $p_{ii}^{(r)}>0$. Por Chapman--Kolmogorov,
\[
p_{jj}^{(m+r+n)}
\ge
p_{ji}^{(n)}\,p_{ii}^{(r)}\,p_{ij}^{(m)}
>0,
\]
luego $m+n+r\in A_j$.

Como también $p_{ii}^{(2r)}>0$, se obtiene igualmente
\[
m+n+2r\in A_j.
\]

Como $d_j$ divide todos los elementos de $A_j$, divide
$m+n+r$ y $m+n+2r$, y por tanto su diferencia:
\[
(m+n+2r)-(m+n+r)=r.
\]

Así, todo $r\in A_i$ es divisible por $d_j$, y por lo tanto
\[
d_j \mid d_i.
\]

\medskip

\textbf{Paso 2: $d_i \mid d_j$.}

Intercambiando los papeles de $i$ y $j$ y repitiendo el argumento,
se obtiene
\[
d_i \mid d_j.
\]

\medskip

Concluimos que $d_i=d_j$.
\end{proof}

\begin{proposition}
Si un estado $i$ en una cadena de Markov tiene período $d_i$, entonces existe un entero $N$ (que depende de $i$) tal que para todo $n \geq N$:
\[
p_{ii}^{( n \cdot d_i)} > 0.
\]

\end{proposition}

Como corolario de la proposición anterior se tiene que si acaso es posible
pasar de $i$ a $j$ en $m$ pasos, entonces también es posible tal transición en
$m + n \cdot d_j$ pasos con $n$ suficientemente grande, suponiendo $d_j \geq 1$.


\section{Primeras Visitas}

Una pregunta fundamental en la teoría de cadenas de Markov es: ¿cuál es la probabilidad de que una cadena alcance un estado (o conjunto de estados) por primera vez exactamente en el paso $n$? Esta pregunta lleva a la definición de \textbf{tiempo de primera visita} y a las correspondientes probabilidades, que son herramientas esenciales para estudiar la estructura y el comportamiento a largo plazo de las cadenas de Markov.

\begin{definition}
Sea $\{X_n : n \in \mathbb{N}\}$ una cadena de Markov con espacio de estados $S$, y sea $B \subseteq S$ un conjunto no vacío de estados. El \textbf{tiempo de primera visita} a $B$ partiendo del estado $i$ se define como:
\[
\tau_B^{(i)} = \min\{n \geq 1 : X_n \in B \mid X_0 = i\},
\]
donde convenimos que $\tau_B^{(i)} = \infty$ si la cadena nunca alcanza $B$ partiendo de $i$.
\end{definition}

\begin{remark}
Para el caso especial donde $B = \{j\}$ es un estado único, escribimos simplemente $\tau_j^{(i)}$ en lugar de $\tau_{\{j\}}^{(i)}$. El evento $\{\tau_j^{(i)} = n\}$ significa que la cadena alcanza el estado $j$ por primera vez en exactamente $n$ pasos, partiendo del estado $i$. Este evento puede escribirse como:
\[
\{\tau_j^{(i)} = n\} = \{X_n = j, X_k \neq j \text{ para todo } k = 1, 2, \ldots, n-1 \mid X_0 = i\}.
\]
\end{remark}

\begin{definition}
Sea $\{X_n : n \in \mathbb{N}\}$ una cadena de Markov con espacio de estados $S$, y sean $i, j \in S$. Definimos la \textbf{probabilidad de primera visita de $i$ a $j$ en $n$ pasos} como:
\[
f_{ij}^{(n)} = \mathbb{P}(\tau_j^{(i)} = n) = \mathbb{P}(X_n = j, X_k \neq j \text{ para } k = 1, \ldots, n-1 \mid X_0 = i).
\]

Alternativamente, es la probabilidad de que, comenzando en $i$, la cadena llegue a $j$ por primera vez exactamente en $n$ pasos.
\end{definition}

\begin{example}[25]
\textbf{Caminata aleatoria sesgada en tres estados.} Considere una cadena de Markov con espacio de estados $S = \{1, 2, 3\}$ y matriz de transición:
\[
P = \begin{pmatrix}
0.3 & 0.7 & 0 \\
0.2 & 0.4 & 0.4 \\
0 & 0 & 1
\end{pmatrix}.
\]

Aquí el estado 3 es absorbente. Queremos calcular $f_{1,3}^{(n)}$, la probabilidad de que, comenzando en el estado 1, la cadena alcance el estado 3 por primera vez en exactamente $n$ pasos.

Para $n = 1$: $f_{1,3}^{(1)} = \mathbb{P}(X_1 = 3 \mid X_0 = 1) = p_{1,3} = 0$ (no hay transición directa desde 1 a 3).

Para $n = 2$: El único camino de 1 a 3 sin pasar por 3 en el paso 1 es $1 \to 2 \to 3$, con probabilidad $p_{1,2} \cdot p_{2,3} = 0.7 \times 0.4 = 0.28$. Por lo tanto, $f_{1,3}^{(2)} = 0.28$.

Para $n = 3$: Hay dos caminos:
\begin{itemize}
    \item $1 \to 1 \to 2 \to 3$: probabilidad $p_{1,1} \cdot p_{1,2} \cdot p_{2,3} = 0.3 \times 0.7 \times 0.4 = 0.084$.
    \item $1 \to 2 \to 2 \to 3$: probabilidad $p_{1,2} \cdot p_{2,2} \cdot p_{2,3} = 0.7 \times 0.4 \times 0.4 = 0.112$.
\end{itemize}

Por lo tanto, $f_{1,3}^{(3)} = 0.084 + 0.112 = 0.196$.
\end{example}

\begin{proposition}
La probabilidad de transición en $n$ pasos desde $i$ a $j$ puede descomponerse en términos de la probabilidad de primera visita. Específicamente, si la cadena aceede de $i$ a $j$ en exactamente $n$ pasos, entonces debe haber alcanzado $j$ por primera vez en algún paso $k \leq n$, y luego hacer la trayectoria completa de $j$ a $j$ en $n - k$ pasos. Formalmente:
\[
p_{ij}^{(n)} = \sum_{k=1}^{n} f_{ij}^{(k)} \cdot p_{jj}^{(n-k)},
\]
donde convenimos que $p_{jj}^{(0)} = 1$ (sin dar pasos, permanecemos en $j$).
\end{proposition}

\begin{proof}
Particionamos todos los caminos de $i$ a $j$ en $n$ pasos de acuerdo al tiempo en el que se realiza la primera visita a $j$:
\[
\{X_n = j \mid X_0 = i\} = \bigcup_{k=1}^{n} \{X_n = j, \tau_j^{(i)} = k \mid X_0 = i\}.
\]

Estos eventos son disjuntos para diferentes valores de $k$. Por lo tanto:
\[
p_{ij}^{(n)} = \mathbb{P}(X_n = j \mid X_0 = i) = \sum_{k=1}^{n} \mathbb{P}(X_n = j, \tau_j^{(i)} = k \mid X_0 = i).
\]

Ahora, si $\tau_j^{(i)} = k$, entonces $X_k = j$ y la cadena no ha visitado $j$ en los pasos $1, \ldots, k-1$. De forma condicional a $X_k = j$, el evento $\{X_n = j\}$ depende únicamente de la evolución de la cadena desde el paso $k$ en adelante, que es independiente del pasado (propiedad de Markov). Así:
\[
\mathbb{P}(X_n = j, \tau_j^{(i)} = k \mid X_0 = i) = \mathbb{P}(\tau_j^{(i)} = k \mid X_0 = i) \cdot \mathbb{P}(X_n = j \mid X_k = j) = f_{ij}^{(k)} \cdot p_{jj}^{(n-k)}.
\]

Sumando sobre $k$, obtenemos la fórmula deseada.
\end{proof}

\begin{definition}
Definimos la \textbf{probabilidad de eventual alcance} de $j$ partiendo de $i$ como:
\[
f_{ij} = \sum_{n=1}^{\infty} f_{ij}^{(n)} = \mathbb{P}(\tau_j^{(i)} < \infty \mid X_0 = i).
\]

Esta es la probabilidad total de que la cadena alcance el estado $j$ en algún momento del futuro, partiendo del estado $i$.
\end{definition}

\begin{remark}
La cantidad $f_{ij}$ representa la probabilidad de que la cadena alguna vez visite el estado $j$ comenzando desde $i$. Tiene interpretaciones muy útiles:
\begin{itemize}
    \item Si $f_{ij} = 1$, decimos que $j$ es \textbf{accesible con certeza} desde $i$: con probabilidad total, la cadena eventualmente alcanzará $j$.
    \item Si $f_{ij} < 1$, existe una probabilidad positiva $1 - f_{ij}$ de que la cadena nunca alcance $j$ partiendo de $i$.
    \item En particular, si $i = j$, entonces $f_{ii}$ es la probabilidad de que la cadena regrese al estado $i$ en algún momento futuro. Esta cantidad es crucial para definir nociones de \textbf{recurrencia} y \textbf{transitoriedad} de los estados.
\end{itemize}
\end{remark}

\begin{example}[26]
\textbf{Caminata aleatoria simétrica: Probabilidad de retorno.} Para la caminata aleatoria simétrica en $\mathbb{Z}$ del Ejemplo 2, calculemos $f_{0,0}$, la probabilidad de que la cadena regrese alguna vez al estado 0 partiendo de 0.

Recordemos que el período es $d_0 = 2$, así que la cadena solo puede retornar al 0 en tiempos pares. Para $n = 2k$ (número par), la probabilidad de retorno en exactamente 2 pasos es:
\[
f_{0,0}^{(2)} = \mathbb{P}(X_2 = 0, X_1 \neq 0 \mid X_0 = 0) = \mathbb{P}(Y_1 = 1, Y_2 = -1) + \mathbb{P}(Y_1 = -1, Y_2 = 1) = \frac{1}{4} + \frac{1}{4} = \frac{1}{2}.
\]

Para el retorno en 4 pasos, debemos evitar retornar en los primeros 2 pasos. Los caminos de retorno en 4 pasos sin tocar 0 intermediamente incluyen secuencias como $(+1, +1, -1, -1)$, $(+1, -1, +1, -1)$, etc. El cálculo directo es más complejo, pero mediante técnicas analíticas (funciones generadoras), puede demostrarse que:
\[
f_{0,0} = \sum_{n=1}^{\infty} f_{0,0}^{(2n)} = 1.
\]

Este es un resultado célebre: la caminata aleatoria simétrica en $\mathbb{Z}$ es \textbf{recurrente}, es decir, con probabilidad 1, la cadena regresa al estado inicial (y de hecho, visita cualquier estado infinitamente a menudo).
\end{example}

\begin{example}[27]
\textbf{Caminata aleatoria sesgada: Transitoriedad.} Consideremos una caminata aleatoria en $\mathbb{Z}$ con $\mathbb{P}(Y_i = 1) = p$ y $\mathbb{P}(Y_i = -1) = q = 1 - p$, donde $p \neq \frac{1}{2}$.

Supongamos $p = 0.6$ (sesgada hacia la derecha). Intuitivamente, si comenzamos en el estado 0, con probabilidad 1 la cadena ``fluye'' hacia $+\infty$, por lo que la probabilidad de retornar al 0 debe ser menor que 1. De hecho, mediante análisis de funciones generadoras, puede demostrarse que:
\[
f_{0,0} = \left(\frac{q}{p}\right) = \frac{0.4}{0.6} = \frac{2}{3} < 1.
\]

Por lo tanto, existe una probabilidad positiva $1 - f_{0,0} = \frac{1}{3}$ de que la cadena nunca retorne al estado 0. En este caso, decimos que la cadena es \textbf{transitoria}: comenzando en cualquier estado, existe una probabilidad positiva de nunca visitarlo nuevamente.
\end{example}

\begin{example}[28]
\textbf{Cadena con estado absorbente.} Reconsideremos el Ejemplo 25 con matriz de transición:
\[
P = \begin{pmatrix}
0.3 & 0.7 & 0 \\
0.2 & 0.4 & 0.4 \\
0 & 0 & 1
\end{pmatrix}.
\]

Ya calculamos que:
\begin{itemize}
    \item $f_{1,3}^{(1)} = 0$ (no hay transición directa desde 1 a 3).
    \item $f_{1,3}^{(2)} = 0.28$ (única ruta $1 \to 2 \to 3$).
    \item $f_{1,3}^{(3)} = 0.196$ (dos rutas).
\end{itemize}

La probabilidad de eventual alcance es $f_{1,3} = \sum_{n=1}^{\infty} f_{1,3}^{(n)}$. Dado que el estado 3 es absorbente y la cadena no puede ``escapar'' a infinito, intuitivamente esperamos que $f_{1,3} = 1$: comenzando en cualquier estado 1 o 2, eventualmente la cadena será absorbida por el estado 3 con certeza.

De hecho, de la matriz podemos ver que desde el estado 1 o 2, la única forma de no alcanzar el estado 3 es permanecer indefinidamente en los estados 1 y 2. Pero estos estados se comunican (viajan entre sí) con probabilidad positiva, y eventualmente la cadena saldrá de este ``grupo'' y será absorbida por 3. Así, $f_{1,3} = 1$.
\end{example}

\vspace{20pt}
\hrule
\vspace{10pt}
{
\color{blue}
\noindent\textbf{--- Termina material del primer parcial ---}
}

\vspace{10pt}
\hrule
\vspace{20pt}


\section{Estados Recurrentes y Transitorios}

Veremos a continuación que los estados de una cadena de Markov pueden ser clasificados, en una primera instancia, en dos tipos, dependiendo de si la cadena es capaz de regresar con certeza al estado de partida. Esta dicotomía es esencial para entender el comportamiento a largo plazo de la cadena: en algunos modelos, la cadena vuelve inevitablemente a ciertos estados; en otros, existe una probabilidad positiva de abandonar un estado para no regresar nunca.

En esta sección utilizaremos la noción de \textbf{primeras visitas} introducida previamente. En particular, para un estado $i \in S$ recordemos que
\[
f_{ii} = \mathbb{P}(\tau_i^{(i)} < \infty \mid X_0 = i)
\]
es la probabilidad de volver a visitar el estado $i$ en algún tiempo futuro (en un tiempo finito), partiendo de $i$.


\begin{definition}[Recurrencia y transitoriedad (I)]
Se dice que un estado $i$ es \textbf{recurrente} si la probabilidad de eventualmente regresar a $i$, partiendo de $i$, es uno; es decir, si
\[
\mathbb{P}(X_n = i \text{ para alguna } n \geq 1 \mid X_0 = i) = 1.
\]
Un estado que no es recurrente se llama \textbf{transitorio}; en tal caso, la probabilidad anterior es estrictamente menor que uno.
\end{definition}

De manera intuitiva, un estado es recurrente si con probabilidad uno la cadena es capaz de regresar eventualmente a ese estado. Más aún, cuando la cadena regresa en algún tiempo finito, por la propiedad de Markov el proceso se \emph{reinicia} desde $i$ y, en ese sentido, puede regresar a $i$ una y otra vez con probabilidad uno. En cambio, el estado se llama transitorio si existe una probabilidad positiva de que la cadena, iniciando en él, ya no regrese nunca a ese estado.

La definición anterior puede enunciarse de manera equivalente en términos de las probabilidades de primera visita.

\begin{definition}[Recurrencia y transitoriedad (II)]
Un estado $i$ es \textbf{recurrente} si $f_{ii} = 1$, es decir, si la probabilidad de regresar a $i$ en un tiempo finito es uno. Análogamente, un estado $i$ es \textbf{transitorio} si $f_{ii} < 1$.
\end{definition}

\begin{remark}
La cantidad $f_{ii}$ se interpreta como una probabilidad de retorno \emph{eventual}. En contraste, las cantidades $p_{ii}^{(n)}$ miden la probabilidad de estar en $i$ \emph{exactamente} después de $n$ pasos. El criterio siguiente conecta ambas familias de cantidades.
\end{remark}

Además de la definición, tenemos el siguiente criterio útil para determinar si un estado es recurrente o transitorio.

\begin{proposition}[Criterio de recurrencia/transitoriedad]
Sea $\{X_n : n \in \mathbb{N}\}$ una cadena de Markov homogénea con espacio de estados $S$ y sea $i \in S$. Entonces:
\begin{enumerate}
    \item[(a)] El estado $i$ es \textbf{recurrente} si, y sólo si,
    \[
    \sum_{n=1}^{\infty} p_{ii}^{(n)} = \infty.
    \]
    \item[(b)] El estado $i$ es \textbf{transitorio} si, y sólo si,
    \[
    \sum_{n=1}^{\infty} p_{ii}^{(n)} < \infty.
    \]
\end{enumerate}
\end{proposition}

\begin{proof}
Fijemos un estado $i \in S$ y supongamos que la cadena inicia en $X_0 = i$. Definamos la variable aleatoria $N_i$ como el número total de visitas al estado $i$ \emph{después} del instante inicial:
\[
N_i = \sum_{n=1}^{\infty} \mathbf{1}_{\{X_n = i\}}.
\]
Por definición, la probabilidad de regresar al menos una vez al estado $i$ en un tiempo finito es $f_{ii}$. Si el proceso regresa, por la propiedad fuerte de Markov evaluada en ese primer instante de retorno, la cadena se reinicia. Por ende, la probabilidad de realizar al menos $k$ retornos, dado que ya presenciamos $k-1$ retornos, sigue siendo $f_{ii}$. Utilizando este argumento de manera inductiva, la suma de los primeros retornos nos da que la probabilidad de que haya al menos $k$ visitas es:
\[
\mathbb{P}(N_i \geq k \mid X_0 = i) = f_{ii}^k, \quad k \geq 1.
\]
Si reescribimos esto en su forma puntual, $\mathbb{P}(N_i = k \mid X_0 = i) = \mathbb{P}(N_i \geq k) - \mathbb{P}(N_i \geq k+1) = f_{ii}^k (1 - f_{ii})$. Es decir, vemos fácilmente que \textbf{la variable aleatoria $N_i$ tiene una distribución geométrica} sobre el conjunto $\{0, 1, 2, \dots\}$ con parámetro de ''éxito'' $1 - f_{ii}$.

Procedamos a calcular $\mathbb{E}[N_i \mid X_0 = i]$ de dos maneras distintas. Usando la definición por variables indicadoras, la linealidad de la esperanza y el teorema de convergencia monótona nos otorgan:
\[
\mathbb{E}[N_i \mid X_0 = i] = \mathbb{E}\left[ \sum_{n=1}^{\infty} \mathbf{1}_{\{X_n = i\}} \;\middle|\; X_0 = i \right] = \sum_{n=1}^{\infty} \mathbb{P}(X_n = i \mid X_0 = i) = \sum_{n=1}^{\infty} p_{ii}^{(n)}.
\]
Alternativamente, como $N_i$ es una variable aleatoria no negativa, su esperanza puede obtenerse sumando las probabilidades de sus colas:
\[
\mathbb{E}[N_i \mid X_0 = i] = \sum_{k=1}^{\infty} \mathbb{P}(N_i \geq k \mid X_0 = i) = \sum_{k=1}^{\infty} f_{ii}^k.
\]

Para concluir, veamos los dos casos posibles:
\begin{itemize}
    \item Si $f_{ii} < 1$ (es decir, $i$ es transitorio), la suma de las colas es una serie geométrica convergente con valor $\frac{f_{ii}}{1 - f_{ii}} < \infty$. Igualando ambos cálculos, encontramos que $\sum_{n=1}^{\infty} p_{ii}^{(n)} < \infty$.
    \item Si $f_{ii} = 1$ (es decir, $i$ es recurrente), los términos en la serie no decaen, por lo que la suma de las colas diverge ($\mathbb{E}[N_i \mid X_0 = i] = \infty$). Por equivalencia, esto implica directamente que $\sum_{n=1}^{\infty} p_{ii}^{(n)} = \infty$.
\end{itemize}
Esto demuestra los incisos (a) y (b) estableciendo el criterio de sumabilidad.
\end{proof}

\begin{remark}
Como se evidenció en la prueba, $1 + \mathbb{E}[N_i \mid X_0 = i] = \sum_{n=0}^{\infty} p_{ii}^{(n)}$ representa el número total esperado de visitas al estado $i$, partiendo de $i$ e incluyendo el tiempo $t=0$. El criterio establece, de manera natural e intuitiva, que un estado es transitorio si su número total esperado de visitas es finito, lo cual concuerda con una probabilidad de retorno estrictamente menor a uno; por el contrario, un estado es recurrente si esperamos visitarlo un número infinito de veces.
\end{remark}

\begin{example}
\textbf{Un estado absorbente es recurrente.} Si $i$ es absorbente, entonces $p_{ii}=1$ y por lo tanto $p_{ii}^{(n)}=1$ para todo $n\geq 1$. En consecuencia,
\[
\sum_{n=1}^{\infty} p_{ii}^{(n)} = \sum_{n=1}^{\infty} 1 = \infty,
\]
de modo que $i$ es recurrente.
\end{example}

\begin{example}
\textbf{Un ejemplo transitorio en dos estados.} Considere la cadena con $S=\{0,1\}$ y
\[
P=\begin{pmatrix}
1-p & p \\
0 & 1
\end{pmatrix}, \quad 0<p<1.
\]
Partiendo de $0$, se tiene $p_{00}^{(n)}=(1-p)^n$, y por lo tanto
\[
\sum_{n=1}^{\infty} p_{00}^{(n)} = \sum_{n=1}^{\infty}(1-p)^n = \frac{1-p}{p} < \infty.
\]
Así, el estado $0$ es transitorio. (En cambio, el estado $1$ es absorbente y por ende recurrente.)
\end{example}






\end{document}
